\documentclass[aspectratio=169, 10pt]{beamer}
\usetheme{metropolis}

% --- Edinburgh colours ----------------------------------------
\definecolor{UoERed}{HTML}{7A2318}
\definecolor{UoEGold}{HTML}{B8860B}
\definecolor{CodeBG}{HTML}{F7F3ED}
\definecolor{CodeFrame}{HTML}{D5C9B1}

\setbeamercolor{palette primary}{bg=UoERed, fg=white}
\setbeamercolor{frametitle}{bg=UoERed, fg=white}
\setbeamercolor{title separator}{fg=UoEGold}
\setbeamercolor{progress bar}{fg=UoEGold, bg=UoEGold!20}
\setbeamercolor{alerted text}{fg=UoERed}
\setbeamercolor{block title}{bg=UoERed!15, fg=UoERed}
\setbeamercolor{block body}{bg=UoERed!5}

% --- Packages -------------------------------------------------
\usepackage[T1]{fontenc}
\usepackage{lmodern}
\usepackage{amsmath, amssymb, mathtools}
\usepackage{listings}
\usepackage{graphicx, booktabs}
\usepackage{tikz}
\usetikzlibrary{arrows.meta, positioning, calc, decorations.pathreplacing}
\usepackage{hyperref}
\hypersetup{colorlinks=true, linkcolor=UoERed, urlcolor=UoERed!70!black}
\setbeamertemplate{navigation symbols}{}

\lstset{
  language=Python,
  basicstyle=\ttfamily\scriptsize,
  keywordstyle=\color{UoERed}\bfseries,
  commentstyle=\color{gray}\itshape,
  stringstyle=\color{UoEGold},
  backgroundcolor=\color{CodeBG},
  frame=single,
  rulecolor=\color{CodeFrame},
  numbers=none, breaklines=true, showstringspaces=false, tabsize=4,
  xleftmargin=4pt, xrightmargin=4pt, aboveskip=6pt, belowskip=4pt,
}

\AtBeginSection[]{
  \begin{frame}
    \vfill\centering
    \usebeamerfont{title}\insertsectionhead\par
    \vfill
  \end{frame}
}

\title{Week 6 --- Constraints and Real-World Models}
\subtitle{Occasionally Binding Constraints \& the Fischer--Burmeister Function}
\author{Dr Juan Zurita}
\institute{Deep Learning for Macroeconomics\\
The University of Edinburgh~$\cdot$~School of Economics}
\date{}

\begin{document}
\maketitle

\begin{frame}{The Constraint Problem}
Many economic models feature \alert{inequality constraints}:

\bigskip
\begin{itemize}
\item Borrowing limits: $a_{t+1} \geq \underline{a}$
\item Non-negativity: $c_t \geq 0$, $k_t \geq 0$
\item Collateral constraints: $b_t \leq \phi\,k_t$
\end{itemize}

\bigskip
These constraints \textbf{bind occasionally} --- sometimes active, sometimes not.

Traditional KKT conditions create a discrete switch that is hard to differentiate.
\end{frame}

\begin{frame}{Fischer--Burmeister Complementarity}
The \textbf{FB function} provides a smooth reformulation:

$$\Phi(a, b) = a + b - \sqrt{a^2 + b^2}$$

\bigskip
\begin{block}{Key property}
$\Phi(a, b) = 0$ if and only if $a \geq 0$, $b \geq 0$, and $ab = 0$.
\end{block}

\bigskip
This is continuously differentiable --- perfect for gradient-based training. It replaces the discrete if/else with a smooth function that autodiff can handle.
\end{frame}

\begin{frame}{Consumption--Savings with Borrowing Constraint}
$\max \sum \beta^t u(c_t)$ s.t.\ $a_{t+1} = (1+r)\,a_t + y_t - c_t$ and $a_{t+1} \geq \underline{a}$

\bigskip
Complementarity: either $a_{t+1} > \underline{a}$ (unconstrained) or $\mu_t > 0$ (constrained).

\bigskip
FB formulation: $\Phi(a_{t+1} - \underline{a},\; \mu_t) = 0$

\bigskip
Add this as an extra residual in the DEQN loss.
\end{frame}

\begin{frame}{Architecture Design}
\begin{itemize}
\item \textbf{Output activation}: softplus for consumption (ensures $c > 0$)
\item \textbf{Width}: 3--4 hidden layers, 64--128 neurons each
\item \textbf{Multiple outputs}: one head for $c$, another for $\mu$ (the multiplier)
\item \textbf{Learning rate schedule}: cosine annealing helps convergence
\end{itemize}

\bigskip
Architecture search is mostly empirical --- start simple, add complexity only if needed.
\end{frame}

\begin{frame}{Summary}
\begin{itemize}
\item Occasionally binding constraints are common and difficult
\item Fischer--Burmeister smooths the complementarity condition
\item DEQNs handle constraints naturally via extra residual terms
\item Architecture choices (activations, depth) matter for convergence
\end{itemize}

\bigskip
\textbf{Next week:} Physics-Informed Neural Networks --- continuous-time economics.
\end{frame}

\end{document}
